\begin{abstract}
This document develops a compact, example-driven path from explicit models of cells and spheres to core computations in singular homology. In Part~1, we construct concrete homeomorphisms—stereographic projection, the open ball–Euclidean space equivalence, and radial maps identifying quotients \(B^n/S^{n-1}\cong S^n\)—and give “standard” boundary-preserving homeomorphisms between simplices, cubes, convex bodies, and balls. In Part~2, we compute \(H_\ast(S^n)\) and \(H_\ast(B^n, S^{n-1})\) via deformation retractions, excision, and long exact sequences of pairs/triples, organize fundamental classes on simplices and spheres, and derive classical consequences: nonexistence of retractions \(B^n\to S^{n-1}\), degree theory on \(S^n\), and a Brouwer-type alternative leading to the fixed-point theorem. Product results are obtained through relative Künneth-style splittings \(H_k(S^n\times Y, P\times Y)\cong H_{k-n}(Y)\), yielding \(H_n(S^n\times S^n)\cong \mathbb{Z}\oplus\mathbb{Z}\). Worked exercises include the homology of punctured Euclidean space, angle-surjectivity for maps without fixed points, and the binomial-rank formula \(H_k(T^n)\cong \mathbb{Z}^{\binom{n}{k}}\). Throughout, the emphasis is on explicit constructions (radial/projection formulas) dovetailed with functorial homological algebra to make global invariants computable from local models.

\textit{Keywords.} Algebraic topology
\end{abstract}
% \\
% \noindent \textit{2020 Mathematics Subject Classification.} Primary 46L55; Secondary 44B20